% Part: first-order-logic
% Chapter: introduction
% Section: soundness-completeness

\documentclass[../../../include/open-logic-section]{subfiles}

\begin{document}

\olfileid{fol}{int}{scp}

\olsection{السلامة والاكتمال}

وسنقدم أيضًا أنساق !!{derivation} لمنطق الرتبة الأولى. وقد طور المناطقة
أنساق !!{derivation} كثيرة، لكنها جميعًا تعرّف علاقة ال!!{derivability}
نفسها بين ال!!{sentence}s. فنقول إن من الممكن \emph{!!{derive}}
~$!A$ من $\Gamma$، ونكتب $\Gamma \Proves !A$، إذا وُجد
!!a{derivation} من نوع
معين مضبوط التعريف. وال!!^{derivation}s دائمًا ترتيبات متناهية من
الرموز؛ فقد تكون قائمة !!{sentence}s، أو بنية أعقد. والغرض من أنساق
ال!!{derivation} أن توفر أداة لتحديد ما إذا كانت !!a{sentence} تلزم عن
مجموعة ما~$\Gamma$. ولكي تؤدي هذا الغرض، يجب أن يكون
$\Gamma \Entails !A$ إذا وفقط إذا كان $\Gamma \Proves !A$.

إذا كان $\Gamma \Proves !A$ ولم يكن $\Gamma \Entails !A$، كان نسق
ال!!{derivation} أقوى من اللازم، لأنه يثبت أكثر مما ينبغي. وتسمى
خاصية أن $\Gamma \Entails !A$ كلما كان $\Gamma \Proves !A$
\emph{السلامة}، وهي شرط أدنى في كل نسق !!{derivation} جيد. ومن جهة
أخرى، إذا كان $\Gamma \Entails !A$ ولم يكن $\Gamma \Proves !A$، كان
نسق ال!!{derivation} أضعف من اللازم، لأنه لا يثبت ما يكفي. وتسمى
خاصية أن $\Gamma \Proves !A$ كلما كان $\Gamma \Entails !A$
\emph{الاكتمال}. وعادة ما يكون إثبات السلامة يسيرًا نسبيًا (بالاستقراء
على بنية ال!!{derivation}s، وهي معرفة استقرائيًا)، أما إثبات الاكتمال
فأصعب.

وللسلامة والاكتمال عدة نتائج مهمة. فإذا أمكن !!{derive} تناقض (مثل
$!A \land \lnot !A$) من مجموعة ال!!{sentence}s~$\Gamma$، سميت
\emph{غير متسقة}. ولا يمكن أن يكون
لمجموعة غير متسقة~$\Gamma$ أي نموذج؛ فهي غير قابلة للإرضاء. ومن الاكتمال
ينتج العكس: فكل $\Gamma$ ليست غير متسقة---أو، كما سنقول،
\emph{متسقة}---لها نموذج. وهذا في الحقيقة مكافئ للاكتمال، وهو
الصورة التي سنبرهن بها الاكتمال فعلًا. وهذه نتيجة عميقة وربما مدهشة:
فمجرد عجزك عن إثبات $!A \land \lnot !A$ من $\Gamma$ يضمن وجود
!!a{structure} تتصف بما تصفها به~$\Gamma$. وبذلك يجيب الاكتمال عن
السؤال: أي مجموعات ال!!{sentence}s لها نماذج؟ الجواب: المجموعات
المتسقة، وهي وحدها.

ولمبرهنتي السلامة والاكتمال نتيجتان مهمتان: مبرهنة التراص ومبرهنة
لوفنهايم--سكولم. وهما نتيجتان مهمتان في نظرية النماذج، ويمكن استعمالهما
لإثبات نتائج كثيرة مثيرة للاهتمام. وقد ذكرنا اثنتين منها: لا يستطيع
منطق الرتبة الأولى التعبير عن أن !!{domain} !!a{structure} متناه، ولا
عن أن مجموعة مجاله !!{nonenumerable}.

وقد استغرق التوصل تاريخيًا إلى كل ذلك وإحكامه وقتًا طويلًا: كيف يُعرّف
نحو منطق الرتبة الأولى ودلالته، وكيف تُعرّف أنساق !!{derivation} جيدة،
وكيف يُبرهن أنها سليمة وكاملة، وكيف يوضّح ما يمكن وما لا يمكن التعبير
عنه في لغات الرتبة الأولى. ونحن نعرف الآن كيف نفعل ذلك، وإن ظل المرور
بجميع التفاصيل مربكًا ومملًا أحيانًا. لكنه مهم أيضًا، لأن المناهج
المطورة هنا للغة منطق الرتبة الأولى الصورية يجري تطبيقها في أنحاء المنطق وعلوم
الحاسوب واللسانيات. ولذلك تؤتي معالجة التفاصيل ثمارها على المدى
البعيد.

\end{document}
